% !TeX root = ../main.tex
\section{Выводы}
\begin{itemize}
    \item \textbf{Характеристики пропускания и селективность фильтра:} Исследованный фильтр нижних частот (ФНЧ) демонстрирует полосу пропускания до 468~МГц по уровню $-1$~дБ и граничную частоту среза 527~МГц по уровню $-3$~дБ. Устройство обладает высокой избирательностью: уровень подавления достигает $-40$~дБ на частоте 657~МГц, а рассчитанный коэффициент прямоугольности составляет 1,25.
    
    \item \textbf{Реактивный характер подавления и КСВН:} В рабочей полосе частот (до 430~МГц) фильтр качественно согласован с волновым сопротивлением тракта 50~Ом, обеспечивая беспрепятственное прохождение мощности в нагрузку (КСВН не превышает $1,6 \dots 2,0$). В полосе заграждения (свыше 550~МГц) КСВН резко возрастает до значений $8,5 \dots 9,3$. Это подтверждает, что подавление сигнала происходит не за счет поглощения мощности (тепловых потерь), а за счет полного отражения электромагнитной волны обратно к источнику.
    
    \item \textbf{Влияние характеристик направленного ответвителя:} Динамический диапазон измерительной установки при анализе отраженных волн ограничен конечной направленностью ответвителя и паразитным просачиванием зондирующего сигнала (уровень шумовой полки/просачивания составляет около $-31,5$~дБм).
    
    \item \textbf{Влияние аттенюаторов в измерительном тракте:} При подключении к тракту ненагруженных аттенюаторов уровень отраженного сигнала падает на величину их удвоенного затухания. Однако из-за упомянутой выше конечной направленности ответвителя, использование аттенюаторов больших номиналов (10~дБ, 20~дБ и выше) приводит к тому, что мощность отраженного сигнала становится сопоставимой с уровнем паразитного просачивания. Это маскирует реальный коэффициент отражения и определяет предел точности измерений в подобных СВЧ-трактах.
\end{itemize}